% Part: incompleteness
% Chapter: arithmetization-syntax
% Section: proofs-in-nd

\documentclass[../../../include/open-logic-section]{subfiles}

\begin{document}

\olfileid{inc}{art}{pnd}
\olsection{自然演绎中的推导}

\begin{explain}
为了将推导算术化，我们需要把推导表示为数字。
由于推导是以公式为节点、每个推理带有一个或两个标签的树，
递归表示是最自然的方法：
我们将一个推导表示为一个元组，其分量包括：
得出最后一个推理之各前提的直接子推导的个数、
这些子推导的表示、末端公式、最后一个推理的解除标签，
以及一个指示最后一个推理类型的数字。
\end{explain}

\begin{defn}
  若 $\delta$ 是自然演绎中的一个推导，则 $\Gn{\delta}$ 归纳定义如下：
  \begin{enumerate}
  \item
    若 $\delta$ 仅由假设~$!A$ 组成，则 $\Gn{\delta}$
    是 $\tuple{0, \Gn{!A}, n}$。其中数~$n$：若该假设未被解除，则取~$0$；否则取其数值标签。
  \item
    若 $\delta$ 以一个有零个、一个、两个或三个前提的推理结束，
    则 $\Gn{\delta}$ 分别为
    \begin{align*}
      & \tuple{0, \Gn{!A}, n, k}, \\
      & \tuple{1, \Gn{\delta_1}, \Gn{!A}, n, k}, \\
      & \tuple{2, \Gn{\delta_1}, \Gn{\delta_2}, \Gn{!A}, n, k}, \text{ 或}\\
      & \tuple{3, \Gn{\delta_1}, \Gn{\delta_2}, \Gn{\delta_3}, \Gn{!A}, n,
        k},
    \end{align*}。
    这里 $\delta_1$、$\delta_2$、$\delta_3$ 是以 $\delta$
    中最后一个推理的各前提为末端的子推导，
    $!A$ 是 $\delta$ 中最后一个推理的结论，
    $n$ 是最后一个推理的解除标签（若该推理不解除任何假设，则 $n=0$），
    而 $k$ 则根据最后一个推理所用的规则，由下表给出：

    \begin{tabular}{lccccccc}
      \text{规则：} & \Intro{\land} & \Elim{\land} & \Intro{\lor} & \Elim{\lor} \\
      $k$： & 1 & 2 & 3 & 4  \\[1.5ex]
      \text{规则：} &  \Intro{\lif} & \Elim{\lif} & \Intro{\lnot} & \Elim{\lnot} \\
      $k$： & 5 & 6 & 7 & 8 \\[1.5ex]
      \text{规则：}  &  \FalseInt & \FalseCl & \Intro{\lforall} & \Elim{\lforall} \\
      $k$： & 9 & 10 & 11 & 12 \\[1.5ex]
      \text{规则：} & \Intro{\lexists} & \Elim{\lexists} & \Intro{\eq} & \Elim{\eq} \\
      $k$： & 13 & 14 & 15 & 16
    \end{tabular}
  \end{enumerate}
\end{defn}

\begin{ex}
  考虑那个非常简单的推导：
  \begin{prooftree}
    \AxiomC{$\Discharge{!A \land !B}{1}$}
    \RightLabel{\Elim{\land}}
    \UnaryInfC{$!A$}
    \DischargeRule{\Intro{\lif}}{1}
    \UnaryInfC{$(!A \land !B) \lif !A$}
  \end{prooftree}
  该假设的Gödel数将是 $d_0 = \tuple{0,
    \Gn{!A \land !B}, 1}$。
  以~$\Elim{\land}$ 的结论为末端的子推导的Gödel数将是 $d_1 = \tuple{1,
     d_0, \Gn{!A}, 0, 2}$（其中 $1$ 是因为 $\Elim{\land}$ 有一个前提，
  接着是结论 $!A$ 的Gödel数，$0$ 表示没有解除任何假设，而 $2$ 是编码 $\Elim{\land}$ 的数字）。
  于是整个推导的Gödel数是
  $\tuple{1, d_1, \Gn{((!A \land !B) \lif
      !A)}, 1, 5}$，即
  \[
  \tuple{1, \tuple{1, \tuple{0, \Gn{(!A \land !B)}, 1}, \Gn{!A}, 0, 2},
    \Gn{((!A \land !B) \lif !A)}, 1, 5}.
  \]
\end{ex}

\begin{explain}
在确定了推导的一种表示方法之后，我们还必须证明可以原始递归地操作此类推导的Gödel数，并表达其本质属性与关系。
有些运算很简单：例如，给定一个推导的Gödel数~$d$，$\fn{EndFmla}(d) = (d)_{(d)_0+1}$
给出其末端公式的Gödel数，
$\fn{DischargeLabel}(d) = (d)_{(d)_0 + 2}$ 给出其解除标签，
而 $\fn{LastRule}(d) = (d)_{(d)_0 + 3}$ 给出指示最后一个推理类型的数字。
有些运算则要困难得多。我们至少会概述如何做到这一点。目标是证明“$\delta$ 是从 $\Gamma$ 到~$!A$ 的推导”这一关系（关于 $\delta$ 和~$!A$ 的Gödel数）是原始递归的。
\end{explain}

\begin{prop}
  下列关系是原始递归的：
  \begin{enumerate}
  \item $!A$ 作为标签为~$n$ 的假设出现在~$\delta$ 中。
  \item $\delta$ 中标签为~$n$ 的所有假设都是~$!A$ 的形式（即，可以在~$\delta$ 中用标签~$n$ 解除假设~$!A$）。
  \end{enumerate}
\end{prop}

\begin{proof}
  我们必须证明，公式的Gödel数与推导的Gödel数之间的相应关系是原始递归的。
  \begin{enumerate}
  \item 我们想证明 $\fn{Assum}(x, d, n)$ 是原始递归的，该关系成立当且仅当 $x$
    是Gödel数为~$d$ 的推导中标签为~$n$ 的某个假设的Gödel数。当Gödel数为 $\tuple{0, x,
      n}$ 的推导是 $d$ 的一个子推导时，便是这种情况。注意，我们对推导的编码方式是
    \olref[cmp][rec][tre]{sec} 中引入的树的编码的一个特例，因此原始递归函数
    $\fn{SubtreeSeq}(d)$ 给出了 $d$ 的所有子推导的Gödel数序列（长度至多为 $d$）。所以我们可以定义
    \[
    \fn{Assum}(x, d, n) \defiff \bexists{i<d}{(\fn{SubtreeSeq}(d))_i =
      \tuple{0, x, n}}.
    \]
  \item 我们想证明 $\fn{Discharge}(x, d, n)$ 是原始递归的，该关系成立当且仅当
    Gödel数为~$d$ 的推导中所有标签为~$n$ 的假设都是Gödel数为~$x$ 的公式。但这个关系成立当且仅当 $\bforall{y<d}{(\fn{Assum}(y, d, n) \lif y
      = x)}$。
  \end{enumerate}
\end{proof}

\begin{prop}
  \ollabel{prop:followsby}
  性质 $\fn{Correct}(d)$ 是原始递归的，它成立当且仅当Gödel数为 $d$ 的推导 $\delta$ 的最后一个推理是正确的。
\end{prop}

\begin{proof}
  此处我们需要证明，对每条推理规则 $R$，关系 $\fn{FollowsBy}_R(d)$ 都是原始递归的，其中 $\fn{FollowsBy}_R(d)$ 成立当且仅当 $d$ 是某个推导 $\delta$ 的Gödel数，且 $\delta$ 的末端公式是由 $\delta$ 的各直接子推导通过对规则 $R$ 的正确应用得出的。

  一个简单的例子是 $\Intro{\land}$ 规则。如果 $\delta$ 以一个正确的 $\Intro{\land}$ 推理结束，它看起来是这样的：
  \begin{prooftree}
    \AxiomC{}
    \RightLabel{$\delta_1$}
    \DeduceC{$!A$}

    \AxiomC{}
    \RightLabel{$\delta_2$}
    \DeduceC{$!B$}

    \RightLabel{\Intro\land}
    \BinaryInfC{$!A \land !B$}
  \end{prooftree}
  那么 $\delta$ 的Gödel数 $d$ 是 $\tuple{2, d_1, d_2,
    \Gn{(!A \land !B)}, 0, k}$，其中 $\fn{EndFmla}(d_1) = \Gn{!A}$，$\fn{EndFmla}(d_2) = \Gn{!B}$，$n=0$，且 $k=1$。因此我们可以将 $\fn{FollowsBy}_{\Intro\land}(d)$ 定义为
  \begin{multline*}
    (d)_0 = 2 \land \fn{DischargeLabel}(d) = 0 \land \fn{LastRule}(d) = 1 \land {}\\
    \fn{EndFmla}(d) = {}\\
    \Gn{(} \concat \fn{EndFmla}((d)_1) \concat \Gn{\land}
    \concat \fn{EndFmla}((d)_2) \concat \Gn{)}.
  \end{multline*}

  另一个简单的例子是 $\Intro\eq$ 规则。该规则没有前提，所以 $(d)_0 = 0$，与假设的情况类似。它也没有解除标签，即 $n=0$。然而，$!A$ 必须具有 $\eq[t][t]$ 的形式，其中 $t$ 是闭项。这里，一个原始递归的定义是
  \begin{multline*}
    (d)_0 = 0 \land \fn{DischargeLabel}(d) = 0 \land {}\\
    \bexists{t<d}{(\fn{ClTerm}(t) \land
      \fn{EndFmla}(d) = {}\\
      \Gn{{\eq}(} \concat t \concat \Gn{,} \concat t \concat \Gn{)})}.
  \end{multline*}

  对于一个更复杂的例子，$\fn{FollowsBy}_{\Intro{\lif}}(d)$ 成立当且仅当 $\delta$ 的末端公式具有 $(!A \lif
  !B)$ 的形式，其中 $\delta_1$ 的末端公式是 $!B$，且 $\delta$ 中任何标签为 $n$ 的假设都具有 $!A$ 的形式。我们可以用原始递归的方式将其表达为
  \begin{multline*}
    (d)_0 = 1 \land {}\\
    \bexists{a<d}{(\fn{Discharge}(a, (d)_1, \fn{DischargeLabel}(d)) \land {}}\\
      \fn{EndFmla}(d) = (\Gn{(} \concat a \concat \Gn{\lif}
      \concat \fn{EndFmla}((d)_1) \concat \Gn{)}))
  \end{multline*}
  （可将 $a$ 视为 $!A$ 的Gödel数）。

  再考虑另一个例子，$\Intro{\lexists}$ 规则。这里，$\delta$ 的最后一个推理是正确的，当且仅当存在公式 $!A$、闭项 $t$ 和变元 $x$，使得 $\Subst{!A}{t}{x}$ 是推导 $\delta_1$ 的末端公式，并且 $\lexists[x][!A]$ 是最后一个推理的结论。所以，$\fn{FollowsBy}_{\Intro{\lexists}}(d)$ 成立当且仅当
  \begin{multline*}
    (d)_0 = 1 \land \fn{DischargeLabel}(d) = 0 \land {} \\
    \bexists{a < d}{\bexists{x<d}{\bexists{t<d}{
          (\fn{ClTerm}(t) \land \fn{Var}(x)  \land {}}}}\\
    \fn{Subst}(a,t,x) = \fn{EndFmla}((d)_1) \land
    \fn{EndFmla}(d) = (\Gn{\lexists} \concat x \concat a)).
  \end{multline*}

  然后我们将 $\fn{Correct}(d)$ 定义为
  \begin{multline*}
    \fn{Sent}(\fn{EndFmla}(d)) \land {}\\
    (\fn{LastRule}(d) = 1 \land
    \fn{FollowsBy}_{\Intro\land}(d)) \lor \dots \lor {}\\
    (\fn{LastRule}(d) = 16 \land \fn{FollowsBy}_{\Elim\eq}(d)) \lor {}\\
    \bexists{n<d}{\bexists{x<d}{(d = \tuple{0, x, n})}}.
  \end{multline*}
  第一行确保 $d$ 的末端公式是一个语句。最后一行涵盖了 $d$ 仅仅是一个假设的情况。
\end{proof}

\begin{prob}
  仿照 \olref[inc][art][pnd]{prop:followsby} 定义以下性质：
  \begin{enumerate}
  \item $\fn{FollowsBy}_{\Elim{\lif}}(d)$，
  \item $\fn{FollowsBy}_{\Elim{\eq}}(d)$，
  \item $\fn{FollowsBy}_{\Elim{\lor}}(d)$，
  \item $\fn{FollowsBy}_{\Intro{\lforall}}(d)$。
  \end{enumerate}
  对于最后一个，你还需要证明可以原始递归地测试Gödel数为 $d$ 的推导的最后一个推理是否满足本征变元条件，即 $\Intro{\lforall}$ 推理的本征变元 $a$ 既不出现在 $d$ 的末端公式中，也不出现在 $d$ 的任何未解除假设中。为此，你可以使用来自
  \olref[inc][art][pnd]{prop:openassum} 的原始递归谓词 $\fn{OpenAssum}$。
\end{prob}

\begin{prop}
  \ollabel{prop:deriv}
  关系 $\fn{Deriv}(d)$ 是原始递归的，它成立当且仅当 $d$ 是某个正确推导 $\delta$ 的Gödel数。
\end{prop}

\begin{proof}
  一个推导 $\delta$ 是正确的，当且仅当它的每一个推理都是对某条规则的正确应用，即当且仅当它的每一个子推导都以一个正确的推理结束。因此，$\fn{Deriv}(d)$ 成立当且仅当
  \[
  \bforall{i<\len{\fn{SubtreeSeq}(d)}}{\fn{Correct}((\fn{SubtreeSeq}(d))_i)}
  \]
\end{proof}

\begin{prop}
  \ollabel{prop:openassum} 关系 $\fn{OpenAssum}(z, d)$ 是原始递归的，它成立当且仅当 $z$ 是Gödel数为 $d$ 的推导 $\delta$ 中某个未解除的假设 $!A$ 的Gödel数。
\end{prop}

\begin{proof}
  如果假设的一次出现在~$\delta$ 的某个子推导中带有标号~$n$，且该子推导结束于一条带有解除标号~$n$ 的规则，则该出现是\emph{解除}的。
  因此，如果 $!A$ 在~$\delta$ 中的至少一次出现未被解除，$!A$ 就是~$\delta$ 的一个\emph{未解除假设}。
  这里必须小心：$\delta$ 可能同时包含 $!A$ 被解除的出现和未被解除的出现。

  考虑序列 $\delta_0$, \dots, $\delta_k$，其中 $\delta_0 = \delta$，$\delta_k$ 是假设 $\Discharge{!A}{n}$（对某个~$n$），且 $\delta_{i+1}$ 是 $\delta_i$ 的一个直接子推导。
  如果存在这样一个序列，其中没有任何 $\delta_i$ 结束于带有解除标号~$n$ 的推理，那么 $!A$ 就是 $\delta$ 的一个未解除假设。

  原始递归函数 $\fn{SubtreeSeq}(d)$ 给出了~$\delta$ 的所有子推导的Gödel编码序列。
  $\delta$ 的任何子推导的Gödel编码序列都是该序列的子序列。
  “作为……的子序列”是一个原始递归关系：$\fn{Subseq}(s, s')$ 成立当且仅当
  $\bforall{i<\len{s}}{\lexists[j<\len{s'}][(s)_i = (s')_j]}$。
  “作为直接子推导”也是原始递归的：$\fn{Subderiv}(d, d')$ 成立当且仅当
  $\bexists{j<(d')_0}{d = (d')_j}$。
  因此我们可以如下定义 $\fn{OpenAssum}(z, d)$：
  \begin{multline*}
    \bexists{s<\fn{SubtreeSeq}(d)}{(\fn{Subseq}(s, \fn{SubtreeSeq}(d))
      \land (s)_0 = d \land {}} \\
    \bexists{n<d}{((s)_{\len{s} \tsub 1} = \tuple{0, z, n} \land {}}\\
      \bforall{i<(\len{s} \tsub 1)}{(\fn{Subderiv}((s)_{i+1}, (s)_i) \land {}}\\
      \fn{DischargeLabel}((s)_i) \neq n))).
  \end{multline*}
\end{proof}

\begin{prop}
  \ollabel{prop:prf-prim-rec}
  设 $\Gamma$ 是一个原始递归的语句集。
  那么，表达“$x$ 是从 $\Gamma$ 中的未解除假设推出 $!A$ 的推导~$\delta$ 的编码，且 $y$ 是 $!A$ 的Gödel编码”的关系 $\Prf[\Gamma](x, y)$ 是原始递归的。
\end{prop}

\begin{proof}
  设“$y \in \Gamma$”由原始递归谓词~$R_\Gamma(y)$ 给出。
  我们需要证明 $\Prf[\Gamma](x, y)$ 是原始递归的。
  该关系成立当且仅当 $y$ 是某个语句~$!A$ 的Gödel编码，且 $x$ 是以 $!A$ 为结束公式、所有未解除假设都在 $\Gamma$ 中的自然演绎推导的编码。

  根据 \olref{prop:deriv}，性质 $\fn{Deriv}(x)$（当且仅当 $x$ 是一个正确的自然演绎推导~$\delta$ 的Gödel编码时成立）是原始递归的。
  因此，我们可以如下定义 $\Prf[\Gamma](x, y)$：
  \begin{align*}
    \Prf[\Gamma](x, y) \defiff {}
    & \fn{Deriv}(x) \land \fn{EndFmla}(x) = y \land {} \\
    & \bforall{z < x}{(\fn{OpenAssum}(z, x) \lif R_\Gamma(z))}.
  \end{align*}
\end{proof}

\end{document}